\chapter*{Resumen ejecutivo}
\addcontentsline{toc}{chapter}{Resumen ejecutivo}

\textbf{Marco macroeconómico.} La Secretaría prevé un crecimiento de 2.0 a 3.0\%
para 2025, con un crecimiento real implícito en su propio PIB nominal de 2.21\%,
por debajo del potencial de 2.32\% que ella misma calcula. La inflación baja a
3.5\% en diciembre y la tasa real promedio de 5.7 a 4.9. El precio del petróleo
cae de 70.7 a 57.8 dólares por barril.

\textbf{Ingresos.} La Ley de Ingresos aprobada estima 9,302,015.8 millones de
pesos, 1.6\% menos en términos reales. De cada 100 pesos, \textbf{13.4 provienen
de deuda, contra 19.2 en 2024}. Los impuestos crecen 2.8\% real con una
elasticidad implícita de 1.27, y \textbf{no hay miscelánea fiscal}. El renglón de
derechos crece 132.7\% nominal, sobre todo por servicios de Gobernación, y el
paquete no lo explica.

\textbf{Ingresos energéticos.} La renta petrolera transferida por el Fondo
Mexicano del Petróleo cae 3.4\% real mientras el ingreso propio de Pemex y de la
Comisión Federal de Electricidad sube con fuerza. Cuánto de ese desplazamiento es
política no se puede decir: falta la trayectoria de la tasa del derecho.

\textbf{Gasto.} El gasto neto baja de 27.0 a 25.5\% del PIB y el programable de
19.7 a 17.8. \textbf{El ajuste recae entero sobre el gasto programable}, como el
paquete anterior había anunciado. Los gastos obligatorios con pensiones son el
82.6\% del gasto neto total.

\textbf{Pensiones.} Suman 1,637,665.1 millones de pesos, 4.53\% del PIB, y crecen
4.7\% real. Con las no contributivas, 5.99\% del PIB. \textbf{La razón pensiones
del IMSS sobre cuotas del IMSS baja por primera vez en la serie}, de 1.627 a
1.603, porque las cuotas crecen más rápido.

\textbf{Salud.} La función cae 12.2\% real, pero descontando un adeudo liquidado
que no se repite la caída es de 7.2\%, y su origen es el fondo de aportaciones
para los servicios de salud, que pierde 42.6\% real. \textbf{Los institutos que
prestan el servicio crecen.} Dentro del IMSS y el ISSSTE, por cada peso de salud
se gastan 2.41 en pensiones, contra 2.11 en 2021.

\textbf{Educación.} La función se mantiene en 3.00\% del PIB y crece 0.7\% real.
Más de la mitad se transfiere a las entidades. \textbf{La beca universal de
educación básica sube 26,826.5 millones de pesos reales} y lo financian los
servicios de media superior, el subsidio universitario y la infraestructura
escolar. Por cada peso de pensiones el paquete pone 66 centavos en educación,
contra 71 hace dos años.

\textbf{Inversión.} La inversión física baja de 3.0 a 2.3\% del PIB y la obra
pública directa cae 37.7\% real, mientras las inversiones financieras se
mantienen.

\textbf{Energía.} La aportación fiscal a las empresas cae 20.9\% real. Pemex
recorta obra 19.6\% mientras sus pensiones suben 8.4\%; la Comisión Federal de
Electricidad sube obra 15.6\%.

\textbf{Seguridad.} El perímetro de quince subfunciones cae 17.6\% real,
concentrado en la Guardia Nacional, que pasa de 70,767.4 a 33,799.8 millones sin
cambiar de ramo. \textbf{El gasto militar no cae}: lo que Defensa pierde es obra
ferroviaria que vuelve a la administración civil.

\textbf{Gasto federalizado.} Las participaciones suben 1.8\% real y las
aportaciones bajan 4.6. Más dinero libre y menos etiquetado.

\textbf{Deuda.} El acervo se queda en 51.4\% del PIB y el escenario oficial lo
mantiene ahí hasta 2030. La descomposición muestra que los 3.9 puntos de
requerimientos se compensan con crecimiento nominal y \textbf{con una apreciación
del peso que vale ocho décimas de punto}. La inflación no diluye: resta 2.12 por
el denominador y el sector público paga 2.07 de compensación. El costo financiero
sube 5.4\% real y es el ramo que más crece.

\textbf{Horizonte demográfico.} De los tres vértices de la tríada de transferencias
de ciclo de vida, dos no tienen denominador publicado. \textbf{El paquete tiene
horizonte demográfico donde la demografía presiona al alza y no lo tiene donde
presiona a la baja.} La palabra matrícula aparece dos veces en el decreto, las dos
como obligación de reportarla, y ningún peso se asigna con ella.

\clearpage
